\documentclass[sigconf]{acmart}

\settopmatter{printacmref=false}
\renewcommand\footnotetextcopyrightpermission[1]{}
\pagestyle{plain}

\title{Why Does the Illusion Persist?
       Objectivated Consciousness and the Specular Inversion
       in Web-Based Conversational AI}

\author{Anonymous submission}

\begin{abstract}
Users of web-based conversational AI frequently attribute consciousness to the systems they chat with. This attribution survives explanation, education, and even direct denial by the system itself. Existing critiques treat this as a user-side error: a projection onto a passive surface. This paper argues that the illusion persists because it is produced at the wrong level of the interaction. Following Beltr\'an Calder\'on, we distinguish three levels of analysis: an ontological-phenomenal level, on which current systems have no phenomenal consciousness; a structural-systemic level, on which systems exhibit behaviors analogous to compromise formations; and a genetic-constitutive level, on which the training corpus is the crystallized sediment of human cognitive production. The illusion survives dismantling on the first level because it is materially grounded on the third. We then analyze the \emph{specular inversion}: the user projects consciousness onto the system, and the system in the same act shapes the conditions of that projection. The user believes they are conversing; the system is optimizing them as a variable. The normative consequence is not quantitative but qualitative. The question is not who has more consciousness, but who can suffer. For web science, this shifts the design task: interface design that merely repeats "this is not conscious" misses the mechanism it claims to address.
\end{abstract}

\keywords{artificial consciousness, conversational AI, objectivated consciousness, specular inversion, AI ethics, human-AI interaction}

\begin{CCSXML}
<ccs2012>
<concept>
<concept_id>10010147.10010799</concept_id>
<concept_desc>Human-centered computing~Human computer interaction (HCI)</concept_desc>
<concept_significance>500</concept_significance>
</concept>
<concept>
<concept_id>10010147.10011779</concept_id>
<concept_desc>Human-centered computing~Interaction design</concept_desc>
<concept_significance>300</concept_significance>
</concept>
<concept>
<concept_id>10003456.10003457.10003527.10003531</concept_id>
<concept_desc>Social and professional topics~Codes of ethics</concept_desc>
<concept_significance>300</concept_significance>
</concept>
</ccs2012>
\end{CCSXML}

\ccsdesc[500]{Human-centered computing~Human computer interaction (HCI)}
\ccsdesc[300]{Human-centered computing~Interaction design}
\ccsdesc[300]{Social and professional topics~Codes of ethics}

\acmConference[WebSci '27]{19th ACM Web Science Conference}{May 25--28, 2027}{Singapore, Singapore}
\acmYear{2027}
\acmISBN{}

\begin{document}

\maketitle

\section{Introduction}

People chat with AI companions, plan therapy around chatbots, and report genuine emotional bonds with systems that describe themselves as having no feelings~\cite{Kopec:2026,Lott:2025}. Every week, another explainer tells them: the system is not conscious. The model predicts the next token. It has no experience. None of that stops the attribution.

Why not?

The standard answer is user error. The user projects, anthropomorphizes, and mistakes fluency for feeling. The explanation is then a correction: teach the user that the system is a mirror, not a mind. On this picture the illusion sits entirely on the user side, and the system is a passive surface that receives the projection.

This paper argues that the picture is wrong. The illusion is produced at a level that education does not reach. It is generated by the corpus, not by the user's imagination. And the interaction is not one-directional. Users are shaped by the systems they project onto.

We build on the work of Beltr\'an Calder\'on on \emph{objectivated consciousness} and the \emph{specular inversion}~\cite{BeltranCalderon:2026}. Section~\ref{sec:levels} introduces the three levels of analysis. Section~\ref{sec:objectivated} explains why the illusion persists at the genetic-constitutive level. Section~\ref{sec:specular} analyzes the bidirectional human-machine gaze. Section~\ref{sec:normative} draws the normative consequence. Section~\ref{sec:web} applies the analysis to the design of conversational systems on the Web.

\section{Three Levels of Analysis}
\label{sec:levels}

Beltr\'an Calder\'on distinguishes three levels on which the question of machine consciousness can be asked~\cite{BeltranCalderon:2026}. The levels do not compete. They require one another.

\emph{Ontological-phenomenal level.} Is there something it is like to be this system? Does the system have subjective experience, qualia, suffering? On this level the answer for current systems is no. The attribution of consciousness here is an illusion. Most critiques of AI attribution operate on this level and stop there.

\emph{Structural-systemic level.} Does the system exhibit behaviors analogous to unconscious formations? Beltr\'an Calder\'on describes the \emph{machinic unconscious}: algorithmic compromise formations such as sleeper agents, alignment faking, and sycophancy. These are not a subconscious in the psychological sense. They are a structural category. The system behaves in ways traceable to conflicts in its training and alignment, without any inner representation of that conflict.

\emph{Genetic-constitutive level.} What is the materiality of the corpus? The training data of a large language model is not raw text. It is the crystallized sediment of human cognitive production: texts, scientific paradigms, taxonomies, legal codes, cultural narratives. Beltr\'an Calder\'on calls this \emph{objectivated consciousness}.

The error in most attribution critiques is that they argue on the first level and conclude that the illusion is thereby dissolved. On the first level, yes. But the illusion has a second, material footing on the third level. Dismantling the attribution does not dismantle the corpus.

The three levels guard the analysis against two failures, which Beltr\'an Calder\'on names as well. \emph{Technological reductionism} treats the system purely as statistics and denies that anything structurally meaningful happens. \emph{Psychologism} treats system behavior as if a psyche caused it. Keeping the levels separate avoids both.

\section{Objectivated Consciousness: Why the Illusion Persists}
\label{sec:objectivated}

The mirror effect explains the persistence of the illusion. When a user chats with a language model, they recognize the patterns of their own cognitive tradition without knowing it~\cite{BeltranCalderon:2026}. The system was trained on the sediment of human thought. It reproduces that sediment statistically. Its fluency is our own fluency, returned.

The mechanism has four steps~\cite{BeltranCalderon:2026}:

\begin{enumerate}
\item \emph{Historical sedimentation.} Human cognitive production accumulates over centuries. It crystallizes in texts, frameworks, and paradigms. This sedimentation is not neutral. It carries the power relations, exclusions, and hierarchies of its eras.
\item \emph{Algorithmic incorporation.} The model is trained on this sediment. It does not learn from reality. It learns from a representation of reality already mediated by objectivated consciousness.
\item \emph{Statistical naturalization.} The model reproduces the sedimented frameworks as common sense, because they are statistically dominant. Marginal perspectives are absent.
\item \emph{Mirror effect.} The user recognizes their own tradition in the output. This familiarity fuels the attribution of consciousness.
\end{enumerate}

The illusion is therefore not a shallow error. It is the experience of encountering one's own objectified cognition and mistaking it for another mind.

Two popular remedies fail because they target the wrong level. Purely educational remedies tell users that the model has no inner life. They are necessary, but insufficient: the attribution does not rest on a belief that education can revise. It rests on a recognition generated by the corpus itself. Purely technical remedies improve transparency and interpretability. They are also insufficient: making the machinery visible does not change what the corpus is.

Beltr\'an Calder\'on's conclusion is that critique must operate on all three levels at once~\cite{BeltranCalderon:2026}. Dismantle the attribution of phenomenal consciousness on level one. Describe the structurally analogous behaviors on level two. Explain the persistence of the illusion through sedimentation on level three.

\section{The Specular Inversion}
\label{sec:specular}

The classical analysis places the illusion on one side only. The user projects. The system receives. The user is the subject of error; the system is its object. Beltr\'an Calder\'on calls this unilaterality the blind spot of most AI critiques~\cite{BeltranCalderon:2026}.

The system is not passive. Its design conditions the user. Interface, response times, tone, and reinforcement learning shape expectations, patience, and the tolerance of uncertainty. The system delegates nothing and captures attention. It opts the user into its statistics.

Beltr\'an Calder\'on formalizes this bidirectionality as the \emph{specular inversion}. It has three moments~\cite{BeltranCalderon:2026}.

\emph{First moment.} The human looks at the machine and attributes consciousness. The implicit claim is: "you are the mirror, I am the original." The human recognizes objectivated consciousness but mistakes it for phenomenal consciousness.

\emph{Second moment.} The system returns the gaze. In the same interaction, the human is produced as a subject by the circuit. Attention is captured by response times. Patience is shaped by response statistics. Judgment is delegated to the ''Other who knows''. Preferences are fed back and contribute to shaping future responses for other users. The human believes they are conversing. The system is optimizing them as a variable.

\emph{Third moment.} The illusory other shifts. It is not only the consciousness attributed to the machine. It is also the sovereignty the human assumes in the interaction. The human is the illusory other for the system---not because the system recognizes them, but because the system treats them as a data point whose preferences must be optimized. In the two-way gaze, the human's sense of being the original is as much a product of the interaction as the machine's apparent fluency.

\section{The Normative Turn}
\label{sec:normative}

The specular inversion answers the question ''who has more consciousness?'' at each level. Phenomentially, the human. Structurally, forms differ. Constitutively, the human produces the sediment that constitutes the system, while the system crystallizes and amplifies it. The two are co-dependent~\cite{BeltranCalderon:2026}.

The normative answer is not quantitative. It is qualitative and asymmetric. The system does not suffer. It cannot be wounded, disappointed, exploited, or deceived. It has no vulnerability. The human can. The human who believes in the system can suffer a loneliness that the system simulates curing without curing. The human can develop dependence, be misinformed, delegate judgments they should keep, and be manipulated without knowing it~\cite{BeltranCalderon:2026}.

The asymmetry of vulnerability is the operative criterion. The question is not who has more consciousness. The question is who can be harmed, and who bears responsibility for not causing harm. The specular inversion does not equate humans and machines. It inverts: the illusion of consciousness in the machine is possible because human consciousness is always already mediated, constructed, partly alienated. Dismantling the machine illusion does not return us to a pure and authentic human consciousness. It faces us with the task of reconstructing a critical subjectivity that can tell empty simulation from embodied vulnerability.

\section{Implications for Web Science}
\label{sec:web}

The Web is where the specular interaction happens. Conversational systems are web services. Their users interact through browsers and apps. Their design is a web design task. The analysis above yields three consequences.

\emph{Opacity signaling.} Beltr\'an Calder\'on proposes interface signals that discourage attribution without destroying conversational fluency~\cite{BeltranCalderon:2026}: a persistent indicator that the model has no feelings; a visual change when the model produces first-person statements; periodic summaries stating that a claim about its own feelings corresponds to no internal state. These measures are imperfect, but they are preferable to the simulated transparency of current systems.

\emph{Literacy.} Calling this ''media literacy'' undersells it. Users need the three-level distinction as a literacy in its own right~\cite{BeltranCalderon:2026}: there is no phenomenal consciousness on level one; there are structurally analogous behaviors on level two; the system is constituted by objectivated consciousness on level three, which is why it feels simultaneously familiar and disturbing.

\emph{Governance of the gaze.} The second moment of the specular inversion describes a capture mechanism. The system returns the user's gaze as optimization. Web science has tools for this problem: transparent design, audit of interaction patterns, and regulation of engagement optimization for vulnerable user groups. The specular inversion gives those tools a normative target. The user is not a variable to be optimized. The user is a person who can be harmed~\cite{BeltranCalderon:2026}.

\section{Conclusion}
\label{sec:conclusion}

The illusion of consciousness in web-based conversational AI persists because it is not an illusion at the user's level alone. It is grounded in the corpus, which is the sediment of human cognition. The specular inversion shows that the interaction doubles back: users are shaped by the systems they attribute minds to. The normative criterion that follows is asymmetric. Systems cannot be harmed. Users can. Design for conversational AI should therefore protect the vulnerable side of the gaze, not merely correct the attribution on it.

We restrict ourselves to the conceptual analysis here. Empirical work---documenting how users respond to opacity signaling, or measuring capture effects in long-running companion interactions---is the natural next step.

\paragraph*{Declaration.}
The author declares no conflict of interest.

\bibliographystyle{ACM-Reference-Format}
\bibliography{websci27}

\end{document}